\section*{Code Availability}\label{sec:code-availability}
The \Argorix compiler, runtime, and offline evidence verifier are openly
available. The evaluation in this paper corresponds to a single, archived
release.

\begin{description}
\item[Repository:] \url{https://github.com/argorixlabs/argorixlang}
\item[Project website:] \url{https://www.argorix-lang.org}
\item[Documentation:] \url{https://www.argorix-lang.org/docs}
\item[Evaluated commit:] \code{6f473f5}
\item[Release version:] v1.0.1
\item[Code license:] Apache License 2.0
\item[Archived DOI:] \url{https://doi.org/10.5281/zenodo.21007563}
\end{description}

The archived DOI resolves to the immutable Zenodo deposit for the release; the
repository URL tracks ongoing development and may diverge from the archived
state.

\section*{Project Website}\label{sec:project-website}
The project website at \url{https://www.argorix-lang.org} hosts the language
overview, and the documentation portal at
\url{https://www.argorix-lang.org/docs} describes language constructs, the
runtime profile, and the evidence schema. Website material is informational and
does not extend the claims evaluated here; the archived repository and Zenodo
deposit remain authoritative for reproduction.

\section*{Artifact Availability}\label{sec:artifact-availability}
The reproducible build pipeline is contained in the \code{paper/} directory of
the archived repository. It comprises the analyzer that normalizes the runtime
snapshot, the PowerShell driver that invokes the offline evidence verifier, the
scripts that render every table and vector figure, the static manuscript checks,
and the pinned LaTeX build. The committed normalized datasets, generated tables,
and generated figures are derived solely from the authorized runtime snapshot
and can be regenerated from it. No executable build tool is committed; the LaTeX
engine is downloaded on demand and validated against a published checksum before
use. Per-request runtime artifacts referenced throughout the paper are the
five-artifact directories described in \cref{sec:evidence}.

\section*{Dataset Availability}\label{sec:dataset-availability}
The empirical study uses an authorized snapshot of 33 runtime request
directories generated by the chatbot runtime. The user authorized publication of
this material, but authorization does not relax data minimization. The committed,
normalized datasets (session-level, event-level, and verifier-outcome records)
contain only aggregated counts and bounded structural fields and may be
published in redacted form. Raw prompt text, secrets, key material, provider
endpoints, absolute filesystem paths, and other sensitive values are not
published. Where the complete corpus is not redistributed, the reason is privacy
and data minimization rather than unavailability of results: the normalized
records are sufficient to recompute every reported total, while the unredacted
inputs remain withheld.

\section*{Reproducibility Instructions}\label{sec:reproducibility-instructions}
The full pipeline normalizes the snapshot, reruns offline bundle verification,
regenerates tables and figures, applies static manuscript checks, and renders
the preprint. The environment and entry points are as follows; build-host values
not fixed by the repository are recorded in the release notes for the archived
deposit.

\begin{description}
\item[Operating system:] developed and validated on Windows; the documented
commands target both PowerShell and POSIX shells.
\item[Rust toolchain:] as pinned by the repository toolchain configuration.
\item[\Argorix version:] v1.0.1 (commit \code{6f473f5}).
\item[Python:] 3.11 or newer, with the pinned \code{paper/requirements.txt}.
\item[Input artifact path:] \code{demo/argorix-chatbot-runtime/generated}
(relative to the repository root).
\item[Output datasets:] \code{paper/data}, with generated tables and figures in
\code{paper/tables} and \code{paper/figures}.
\end{description}

The canonical entry point regenerates all artifacts:

\begin{lstlisting}
python -m pip install -r paper/requirements.txt
powershell -NoProfile -ExecutionPolicy Bypass \
  -File paper/scripts/build_paper.ps1 -Target all
\end{lstlisting}

The offline evidence verifier can be rerun in isolation over the snapshot:

\begin{lstlisting}
powershell -File paper/scripts/verify_runtime.ps1 \
  -InputRoot demo/argorix-chatbot-runtime/generated \
  -OutputPath paper/data/verification-results.json
\end{lstlisting}

The expected outcome is deterministic: 27 of 33 complete EvidenceBundles verify
(27/27 of the complete set), and the six source-only directories are retained as
unavailable rather than discarded. Successful verification confirms offline
artifact integrity only; it is not policy approval, a security proof, or
certification.

\section*{Required Artifact Changes Before Expanded Claims}\label{sec:required-changes}
The paper's stronger v0.2 claims require implementation evidence, not prose
alone.  Required code and artifact changes are: implement the typed policy
decision lattice in \cref{tab:policy-lattice}; add known rules for executable
provider allowlists, external adapters, network denial, secret and key-material
denial, passport residency, source digest presence, bundle digest matching, and
report consistency; add a source digest to the EvidenceBundle and bind it to the
compiled bytecode; generate the controlled cases in
\cref{tab:controlled-matrix}; implement tamper tests for modified bytecode,
trace, security report, and source; record side-effect execution status per
case; and publish latency, evidence-size, trace-overhead, compile-time, and
policy-evaluation metrics for the ablations in \cref{tab:ablation}.  Until
those artifacts exist, the current paper must report them as Required
experiment, To be generated, or Do not claim until measured.

\section*{Reviewer Checklist}\label{sec:reviewer-checklist}
A reviewer should be able to answer yes to the following bounded claims and no
to the corresponding overclaims.  The paper claims a compiled governance
language/runtime, a fail-closed provider boundary, and offline internal evidence
consistency for 27 complete bundles.  It does not claim security certification,
real identity verification, cryptographic attestation, live ANS/DNS federation,
production sandboxing, broad prompt-injection defense, or a large empirical
security corpus.  Counts over the repeated snapshot are descriptive.  Unknown
policy rules are treated as a negative diagnostic finding.  Source integrity,
signatures, append-only logs, workload diversity, and controlled ablations are
future assurance requirements unless their artifacts are generated.

\section*{Ethics and Privacy Statement}\label{sec:ethics}
This study analyzes only material that the user authorized for publication, and
it applies data minimization throughout. We do not reconstruct prompts, and we
publish no secrets, credentials, endpoints, or otherwise sensitive content. The
analysis inspects only the bounded \code{trace.injected.content} schema field,
which is present in 0 of 27 complete, valid traces; the six source-only
directories contain no trace and are not assessable. Because no prompt content is
available, we perform no prompt-level qualitative analysis, quote no prompts, and
report no prompt-injection success rate. The absence of prompt text reflects a
deliberate minimization decision and the studied runtime configuration; it is
not evidence of the absence of risk. The work involves no human subjects,
personal data collection, or interaction beyond the authorized artifact
snapshot.

\section*{Author Contributions}\label{sec:contributions}
The authors jointly contributed to the conception of the system, the language
and runtime design, the implementation under evaluation, the observational
analysis, and the writing and revision of the manuscript. Detailed CRediT role
assignments will be finalized for the camera-ready version.

\section*{Funding}\label{sec:funding}
No dedicated external grant funded the preparation of this preprint. Any
institutional support will be acknowledged in the camera-ready version.

\section*{Competing Interests}\label{sec:competing-interests}
The authors are affiliated with the organizations developing \Argorix and
disclose this relationship as a potential competing interest. The authors
declare no other competing financial or non-financial interests.

\section*{License}\label{sec:license}
The \Argorix implementation is released under the Apache License 2.0. The
archived software deposit is citable through its Zenodo DOI. The license under
which this manuscript text is distributed will be confirmed at submission.
